\section{Implementation Contract}

The implementation contract has three parts: validity-preserving crossover and
mutation operators, deterministic fixed-seed frontier metadata, and selected
frontier export. Operators may return valid children, explicit valid fallbacks,
or rejected-child records; invalid silent outputs falsify the contract.

The atlas view makes the frontier decision explicit. Offspring creation,
repair/fallback, frontier observation, and selected-index export are different
events. A selected entry need not be best in every objective; it is a chosen
trade-off point whose generation, objective vector, validity status, and
selection index must travel with the RPLAN/RCTX package.

The implementation's frontier object makes the record concrete. Each frontier
entry carries a one-based district assignment, objective values for edge cut,
seats, and VRA deficit, a non-dominated flag, generation found, plan-level
validity status, and optional audit-certificate identity fields for exported
entries. The final NDJSON line records population size, generation count,
frontier size, base seed, runtime, output version, and a SHA-256 hash of the
plan lines. This is the evidence that a row was observed in a particular search
run, not evidence that no better row exists outside that budget.

A selected-frontier example is deliberately small. Suppose index 1 has edge cut
137 and county splits 8 while index 0 has edge cut 130 and county splits 13.
Index 1 is not globally ``better''; it is the selected trade-off. Exporting it
therefore records selected index 1, frontier size, base seed, generation,
objective vector, and validity status before writing the selected RPLAN/RCTX
package.

\begin{table}[htbp]
\centering
\small
\begin{tabular}{p{0.22\linewidth}p{0.34\linewidth}p{0.29\linewidth}}
\toprule
Layer & Record & Claim boundary \\
\midrule
Operator-valid & parent ids, seed, parameters, validity status & child validity \\
Operator-fallback & parent ids, failed operator, fallback reason & validity preservation \\
Frontier-recorded & objective vectors and dominance metadata & search result under budget \\
Selected-packaged & chosen entry and audit package id & auditable final plan \\
Selection-deferred & frontier entry without selection record & comparison only \\
\bottomrule
\end{tabular}
\caption{Evolutionary comparison record layers.}
\end{table}

\begin{table}[htbp]
\centering
\small
\begin{tabular}{p{0.27\linewidth}p{0.29\linewidth}p{0.30\linewidth}}
\toprule
Artifact field & Current value family & Interpretation \\
\midrule
Frontier row & plan, objectives, generation, validity & Observed comparison entry \\
Metadata row & population, generations, base seed, hash & Reproducible run envelope \\
Selected package & selected index, frontier size, RPLAN/RCTX paths & Chosen entry exported for audit \\
Lineage extra & method, seed, generation, objectives, validity & Bridge from package back to frontier \\
Audit certificate & profile, context hash, result, content hash & Final-plan validity evidence \\
\bottomrule
\end{tabular}
\caption{Concrete U.19 fields that separate frontier evidence from audit evidence.}
\end{table}

The implementation therefore separates producing candidates, observing a
frontier, and selecting an entry for downstream audit. Frontier membership is a
search-output claim. Selected-package status is an audit-readiness claim.
Selection-deferred status means the entry can be compared but should not be
reported as a chosen plan without the U.14 selection record.
